\documentclass[10pt]{article}
\usepackage[margin=20mm]{geometry}
\usepackage{amsmath,amssymb,mathtools,fontspec,unicode-math,xurl,hyperref}
\IfFontExistsTF{DejaVu Serif}{\setmainfont{DejaVu Serif}}{\setmainfont{Latin Modern Roman}}
\setmathfont{latinmodern-math.otf}
\setlength{\parindent}{0pt}\setlength{\parskip}{6pt}\sloppy
\title{Scientific initial review: M25-0002
Galois sections and \(p\)-adic period mappings}
\author{}
\date{}
\begin{document}
\maketitle
Abstract: The reviewed object is the 59-page author manuscript arXiv:2204.13674v1, dated 28 April 2022, by L. Alexander Betts and Jakob Stix. One demonstrated defect concerns the period-map comparison used in the core finiteness argument; one concerns the unnecessary generality of the abstract principal trichotomy. Both admit explicit counterexamples to the particular printed assertions. Neither counterexample disproves the main finiteness theorem. The 2025 journal Version of Record was not acquired or compared, and these findings are not attributed to it.\par
\section*{Reviewed object and scope}
The frozen journal identity is Annals of Mathematics 201 (2025), 79–166, DOI 10.4007/annals.2025.201.1.2 [S2]. The publisher records revision on 24 May 2024 and acceptance on 28 May 2024, substantially after the reviewed v1 [S1]. The arXiv record exposed in the version check lists only v1. The author’s publication list links the same arXiv identifier and journal article [S3]. Bounded searches of the exact title and identifier with correction, erratum, and revision terms did not identify a specific separate correction. This is not a proof that none exists.\par
All 59 PDF pages, including references and Appendix A, were read. The active author TeX was inspected statically at the definitions and proof passages relevant to the findings. The input sequence in `main.tex`, lines 133–143, excludes the older file `prelims.tex`; material in that inactive file is not evidence of a defect in the PDF. Source-manifest file bindings were verified by size and SHA-256. Source pages 16, 33, 40–42, 44, 46–47, 49, and 51 were also visually inspected. No author code was executed.\par
The main assertion concerns smooth projective curves of genus at least two over number fields without a CM subfield, with finiteness of the local image of Selmer sections. The intermediate local theorem uses a self-conjugate place. Smoothness, projectivity, the genus restriction, the place condition, and the Hodge–Tate weight restrictions are not discarded in this review. Reading the whole manuscript does not amount to independently reproving every cited theorem of \(p\)-adic Hodge theory.\par
\section*{F01: the formal period map uses transport in the wrong direction}
Classification: Major; demonstrated invalid core derivation.\par
\subsection*{Location and assertion}
Definition 3.1, continued on p. 16, defines the period filtration by pulling back along horizontal transport \[T^\nabla:\mathcal O_U\otimes_{K_v}\mathcal E_{y_0}\longrightarrow \mathcal E|_U.\] In Section 5.2, pp. 40–42, the connection is written \(\nabla=d+\omega\), the filtration is made constant in a chosen frame, and equation (5.1) specifies \(dT=-\omega T\), \(T(0)=I\). The manuscript then constructs \(\widehat\Phi\) from the first columns of \(T\), identifies these with the analytic period coordinates, and uses this identification in Lemmas 5.5 and 5.6. The exact source locations are `big\_monodromy.tex`, lines 76–112 and 116–137. Definition 3.1 is `reps\_in\_families.tex`, lines 10–29.\par
With the displayed direction of \(T^\nabla\), the pullback of a subspace \(F\) is \(T^{-1}F\). Its coordinates in the reference fibre are therefore columns of \(T^{-1}\), not columns of \(T\). This distinction can change the Zariski closure, so it cannot be dismissed merely as a different parametrization of the same image.\par
\subsection*{A counterexample satisfying the stated setup}
Take \(K_v=\mathbb Q_3\), \(Y=\mathbb A^1_{\mathbb Q_3}\), \(y_0=0\), \(\mathcal E=\mathcal O_Y^3\), and the constant filtration \[F^i\mathcal E=
\begin{cases}
\mathcal E,&i\leq0,\\
\mathcal O_Y e_1,&i=1,\\
0,&i\geq2.
\end{cases}\] It is exhaustive and separated, and all graded pieces are vector bundles. Set \[T=\begin{pmatrix}1&0&0\\t&1&0\\t^2&t&1\end{pmatrix},
\qquad
T^{-1}=\begin{pmatrix}1&0&0\\-t&1&0\\0&-t&1\end{pmatrix},
\qquad
\omega=\begin{pmatrix}0&0&0\\-1&0&0\\-t&-1&0\end{pmatrix}dt .\] Direct multiplication gives \(TT^{-1}=T^{-1}T=I\), \(T(0)=I\), and \[\omega=-T'T^{-1}\,dt,\qquad dT+\omega T=0.\] The algebraic connection is flat: on this smooth one-dimensional base, both \(d\omega\) and \(\omega\wedge\omega\) vanish. Its horizontal frame is the polynomial matrix \(T\), invertible on the entire affine line. In particular, on the closed unit disc \(U=\operatorname{Sp}(\mathbb Q_3\langle t\rangle)\) it has a full horizontal basis and \(H^0_{\mathrm{dR}}(U/\mathbb Q_3)=\mathbb Q_3\). Even Griffiths transversality, which Definition 3.1 does not require, holds, since \(\nabla(F^1)\subset F^0\otimes\Omega^1\).\par
The actual period map prescribed by Definition 3.1 is \[t\longmapsto T(t)^{-1}\langle e_1\rangle=[1:-t:0].\] Its Zariski closure in \(\mathbb P^2\) is the line \(X_2=0\): a polynomial in the remaining affine coordinate vanishing on all \(t\in\mathbb Q_3\) with \(|t|_3\leq1\) is zero. The printed formal construction instead gives \[\widehat\Phi(t)=[1:t:t^2],\] whose scheme-theoretic image is the conic \(X_0X_2-X_1^2=0\). Indeed, on \(X_0\ne0\), the kernel of \(\mathbb Q_3[u,v]\longrightarrow\mathbb Q_3[[t]]\), \((u,v)\mapsto(t,t^2)\), is exactly \((v-u^2)\): divide by the monic polynomial \(v-u^2\) and use injectivity of \(\mathbb Q_3[u]\to\mathbb Q_3[[t]]\). For the actual map the kernel is \((v)\). These are different ideals. Shrinking the disc does not remove the difference. The same matrices over \(\mathbb C\) give the same discrepancy for Lemma 5.6.\par
The independently written checking script performs exact polynomial matrix arithmetic; both inverse products and the differential-equation residual were checked coefficient by coefficient. The ideal arguments above are algebraic proofs, not conclusions drawn from sampling.\par
\subsection*{Counterevidence, dependency, and repair}
Remark 5.7 on p. 42 explicitly returns to the correct intrinsic description by pullback along \(\widehat T\). This supports the intended construction, but contradicts the preceding column formulas rather than validating them. Changing the direction of transport only in the interpretation would conflict with Definition 3.1 and equation (5.1).\par
The dependency is direct: Lemmas 5.5 and 5.6 are used immediately to prove Proposition 5.3. That comparison enters Lemma 5.9 and Corollary 5.10, and hence the finiteness applications in Corollaries 6.10 and 6.16, Theorem 6.5, and Theorem A together with Theorem 7.1. Under the prescribed review standard, a demonstrated invalid derivation on this core path is Major even when a local repair is available.\par
The example is a filtered flat bundle allowed by the general comparison setup; it is not asserted to arise from an abelian-by-finite family. It does not refute Theorem A or the true comparison of analytic period images. A uniform repair of the displayed construction is to use columns of \(T^{-1}\) on both analytic sides and in \(\widehat\Phi\). Since \(T(0)=I\), the inverse exists formally and on a sufficiently small analytic neighbourhood. Inversion commutes with the field embedding \(\iota:K_v\hookrightarrow\mathbb C\), so the corrected complex coordinates are the images of the corrected \(v\)-adic coordinates. The kernel base-change argument then applies to these coordinates. Thus the defect and its algebraic repair are both established here; the report does not claim that the principal finiteness conclusion is false.\par
\section*{F02: the principal trichotomy allows algebras for which its count fails}
Classification: Minor relative to the paper’s core result; demonstrated false extra generality.\par
\subsection*{Location and complete example}
Proposition 6.7 on p. 46, also summarized as Proposition 1.8, permits the symplectic pairs of Definitions 4.3–4.4 (pp. 33–34). Those definitions allow a finite-dimensional \(\mathbb Q_p\)-algebra \(A\) without requiring it to be split or reduced. The proposition sets \(\Sigma=\operatorname{Hom}_{\mathbb Q_p\text{-alg}}(A,\mathbb Q_p)\). Its alternatives (a) and (b) require an orbit of local size at least four, whereas (c) requires at least \(\dim_{\mathbb Q_p}(A)/(d+1)\) elements in local orbits of size less than four. See `locus.tex`, lines 142–163.\par
Let \(K=\mathbb Q\), \(p=3\), \(v=3\), and \[A=\mathbb Q_3[\epsilon]/(\epsilon^3)\] with trivial \(G_\mathbb Q\)-action. Let \(E/\mathbb Q\) be the elliptic curve \(y^2=x^3-x\), let \(B=H^1_{\mathrm{\acute et}}(E_{\overline{\mathbb Q}},\mathbb Q_3)\), and put \[V=A\otimes_{\mathbb Q_3}B,\qquad
L=A\otimes_{\mathbb Q_3}\mathbb Q_3(-1).\] The cup-product pairing on \(B\) extends to a perfect alternating \(A\)-bilinear, \(G_\mathbb Q\)-equivariant pairing \(V\times V\to L\). Here \(V\) is free of rank \(2\) over \(A\), so \(d=1\), and \(L\) is free of rank one. The actions are continuous and compatible with all structures.\par
The place is self-conjugate because \(\mathbb Q\) has no CM subfield. The elliptic curve has discriminant \(64\), hence good reduction outside \(2\). Taking \(S=\{2,3,\infty\}\), \(B\) is unramified, pure and integral of weight one away from \(S\), is de Rham at \(3\), and has Hodge–Tate weights \(\{0,1\}\), as for the smooth proper cohomology used throughout the manuscript. As a \(\mathbb Q_3\)-representation \(V\) is three copies of \(B\), so it satisfies the same hypotheses; \(A\) is unramified everywhere. These standard cohomological facts, not a finite numerical test, supply the arithmetic hypotheses of the example.\par
Every \(\mathbb Q_3\)-algebra homomorphism \(A\to\mathbb Q_3\) sends \(\epsilon\) to zero, since a field has no nonzero nilpotents. Thus \(\Sigma\) has exactly one element, fixed by \(G_\mathbb Q\) and \(G_v\). Alternatives (a) and (b) are impossible because its local orbit has size \(1<4\). Alternative (c) would require \[1\ \geq\ \frac{\dim_{\mathbb Q_3}A}{d+1}=\frac32,\] which is false. All three alternatives therefore fail.\par
\subsection*{Failed inference and restricted valid argument}
In the preliminary reduction on p. 46, the proof argues from an orbit proportion bounded below by \(1/(d+1)\) and \(\#\Sigma\leq\dim_{\mathbb Q_p}A\) to (c). The inequality has the wrong direction for this inference. In the example it compares \(1\) with \(3\). The natural map \(A\to\prod_{\psi\in\Sigma}\mathbb Q_p\) forgets its nilpotent directions, which contribute to the denominator but not the count.\par
There is decisive counterevidence against treating this as a failure of the main theorem: Remark 6.3 on p. 44 states that an interpolating \(A\) in the Lawrence–Venkatesh locus is isomorphic as an algebra to \(\mathbb Q_p^\delta\). Remark 6.6 on p. 46 displays this fact for geometric degree-zero cohomology. Thus the actual application already has \(\#\Sigma=\dim_{\mathbb Q_p}A\).\par
Restricting Proposition 6.7 to split reduced algebras \(A=\mathbb Q_p^\delta\) removes this counterexample and validates the orbit reduction. Merely requiring a finite étale algebra does not ensure this equality, because the homomorphisms in the statement have target \(\mathbb Q_p\), not \(\overline{\mathbb Q}_p\). For a transitive split component, write \(m=\dim W\), \(s=\#\Sigma_{<4}\) and \(\delta=\#\Sigma\). Failure of (b), the integral Hodge–Tate dimensions, and the averaged weight identity give \[\frac{m\delta}{2}
 \leq\frac{m-1}{2}(\delta-s)+ms,
\qquad
\delta\leq(m+1)s\leq(d+1)s.\] The last step uses isotropy, \(m\leq d\). Summing the transitive split components yields the stated count. This reconstruction confirms why the restriction is sufficient for the intended application. The false abstract generality deserves correction, but no core failure from this second issue has been demonstrated.\par
\section*{Independent positive checks and boundaries}
\subsection*{Four Lagrangians and the degree-four threshold}
The argument in Sublemma 6.18 on p. 51 was reconstructed. Over an algebraic closure choose complementary Lagrangians \(E_0=\langle e_i\rangle\) and \(F_0=\langle f_i\rangle\), and graphs of maps \(\psi_3(e_i)=f_i\) and \(\psi_4(e_i)=\lambda_i f_i\), where the nonzero \(\lambda_i\) are distinct. All pairs are complementary except possibly the two graphs. If a nonzero isotropic \(W\) met each Lagrangian in dimension at least \(\dim(W)/2\), complementary pairs force \(\dim W\) to be even and \[W=(W\cap E_0)\oplus(W\cap F_0),\qquad
\psi_3(W\cap E_0)=\psi_4(W\cap E_0)=W\cap F_0.\] Hence \(W\cap E_0\) is invariant under the diagonal map \(\psi_3^{-1}\psi_4\). A nonzero invariant subspace for a diagonal map with distinct eigenvalues contains an eigenvector, by polynomial spectral projection. Thus \(e_i\) and \(f_i\) both belong to \(W\) for some \(i\), contradicting \(\omega(e_i,f_i)=1\). The proper incidence projection then gives the required proper closed bad locus; \(r>4\) follows from the first four factors.\par
As a separate finite check, the reviewer enumerated all subspaces of \(\mathbb F_5^4\) of dimensions one and two using reduced row-echelon matrices: 156 lines and 806 planes, of which 156 planes are isotropic. None of the 312 nonzero isotropic subspaces meets all four chosen Lagrangians in at least half its dimension. This finite instance is supporting evidence; the preceding argument supplies the universal proof. The strict dimension estimate at the degree-four threshold also checks exactly: \[4\,\frac{d(d+1)}2-d(2d+1)=d>0\quad(d>0).\]\par
\subsection*{Arithmetic bounds, local conditions, and dependencies}
The final inequality used in Proposition 6.20 was checked algebraically. Writing \(q-1=2k\), with \(g\geq2\) and \(k\geq2\), \[4k^{g-1}\geq 2^g k\geq2gk>(2g-1)k+1.\] Together with \(\zeta(2g)<9/8\) and \(2k\geq9n_v\), this gives the required upper bound for the proportion of short orbits. Exact rational tests for 741 pairs \((g,k)\) supplement this argument. For the stated genus three illustrations, \(q=11\) gives \(d_q=25\) and \(9/500<1/26\); \(q=23\) gives \(d_q=55\) and \(9/2662<1/56\). The Chinese-remainder congruences used with Dirichlet’s theorem are compatible and prime to their modulus. These checks concern the inequalities, not practical runtime of a point-finding algorithm.\par
The determinant and induction argument for averaged Hodge–Tate weights was checked with the local degree factors retained. For Proposition 2.16 the required endomorphism dimension bound can be justified by Frobenius descent: a shortest normalized linear relation among Frobenius-fixed endomorphisms, after applying Frobenius and subtracting, forces all coefficients to be Frobenius-fixed, hence in \(\mathbb Q_p\). Averaging only under \(G_v\) requires care when its unramified fixed field exceeds \(\mathbb Q_p\); this does not provide a counterexample to the dimension bound itself.\par
The closed-polydisc condition in the analytic finiteness step, the properness used for the global finite cover, the self-conjugate-place restriction, and the basepoint/conjugacy compatibility of the section construction were checked along their displayed applications. The degree-two and degree-three unramified extensions used when the residue field has two elements give the necessary coprime subgroup indices. The manuscript explicitly includes the Scholze erratum in its filtered-derived discussion; it was not silently treated as the uncorrected source.\par
\section*{Assessment and remaining limits}
The supported result is one Major finding and one Minor finding for the specified v1. The first invalidates a printed core proof step, with a concrete correction and no counterexample to the final finiteness conclusion. The second concerns generality excluded by the actual application. They are not combined into an unsupported claim that Theorem A is false.\par
The review covers the entire manuscript at reading level and the displayed proof chains at the level described above. The full external proofs of the Scholze–Shimizu comparison results, the Faltings finiteness input, and the Lawrence–Venkatesh monodromy construction were not all independently rederived. No journal-final comparison was performed. The report therefore neither certifies global correctness nor establishes which issues, if any, survived the 2024 journal revision.\par
\section*{Compact conclusion table}
Item | Judgment | Repair and remaining decision\par
F01, period transport | Major; core displayed step invalid | Use the inverse transport matrix consistently. The explicit connection refutes the printed step; it does not disprove Theorem A. The corrected full proof remains to be adopted and checked.\par
F02, Proposition 6.7 | Minor; nonreduced generality | Exclude the nilpotent-algebra counterexample by the reduced/split hypothesis available in the application. No core dependency failure shown there.\par
Scope and decision | v1 initial review complete | Full journal-final comparison and all external comparison/monodromy proofs remain outside the established coverage. Author adoption was not requested; editorial disposition remains pending.\par
S3 L. A. Betts and J. Stix, Galois sections and \(p\)-adic period mappings, arXiv:2204.13674v1 (28 April 2022), reviewed PDF and bound author sources. https://arxiv.org/abs/2204.13674v1. Annals of Mathematics, article metadata for Galois sections and \(p\)-adic period mappings, 201 (2025), 79–166, DOI 10.4007/annals.2025.201.1.2. https://annals.math.princeton.edu/2025/201-1/p02. J. Stix, author publication list, entry linking the 2025 article and arXiv:2204.13674. https://www.math.uni-frankfurt.de/\textasciitilde{}stix/publikationen.html.\par
\end{document}
